\newcommand{\titulus}{\nomenFesti{S. Irenæi, Episcopi, Martyris.}
\dies{Die 28. Iunii.}}
\newcommand{\sineobmv}{Sine officium B.M.V.}
\newcommand{\oratio}{\pars{Oratio.}

\noindent Deus, qui beáto Irenǽo, epíscopo, tribuísti, ut veritátem doctrínæ pacémque Ecclésiæ felíciter confirmáret, concéde, quǽsumus, eius intercessióne, ut nos, fide et caritáte renováti, ad unitátem concordiámque fovéndam semper simus inténti.

\pars{Pro pace in universo mundo.} \scriptura{Sir. 50, 25; 2 Esdr. 4, 20; \textbf{H416}}

\vspace{-4mm}

\antiphona{II D}{temporalia/ant-dapacemdomine.gtex}

\vfill

\noindent Deus, a quo sancta desidéria, recta consília et iusta sunt ópera: da servis tuis illam, quam mundus dare non potest, pacem; ut et corda nostra mandátis tuis dédita, et hóstium subláta formídine, témpora sint tua protectióne tranquílla.

\noindent Per Dóminum nostrum Iesum Christum, Fílium tuum, qui tecum vivit et regnat in unitáte Spíritus Sancti, Deus, per ómnia sǽcula sæculórum.

\noindent \Rbardot{} Amen.}
\newcommand{\invitatorium}{\pars{Invitatorium.}

\vspace{-4mm}

\antiphona{E}{temporalia/inv-regemmartyrumsimplex.gtex}}
\newcommand{\hymnusmatutinum}{\pars{Hymnus}

\cuminitiali{I}{temporalia/hym-BeateMartyr.gtex}}
\newcommand{\matversus}{\noindent \Vbardot{} Iustum dedúxit Dóminus per vias rectas.

\noindent \Rbardot{} Et osténdit illi regnum Dei.}
\newcommand{\lectioi}{\pars{Lectio I.} \scriptura{1 Sam. 26, 5-12}

\noindent De libro primo Samuélis.

\noindent In diébus illis: Surréxit David et venit ad locum, ubi erat Saul. Cumque vidísset locum, in quo dormiébat Saul et Abner fílius Ner princeps milítiæ eius, Saulem dormiéntem in carrágine et réliquum vulgus per circúitum eius, ait David ad Achímelech Hetthǽum et Abísai fílium Sárviæ fratrem Ioab dicens: «Quis descéndet mecum ad Saul in castra?». Dixítque Abísai: «Ego descéndam tecum».

\noindent Venérunt ergo David et Abísai ad pópulum nocte et invenérunt Saul iacéntem et dormiéntem in carrágine et hastam fixam in terra ad caput eius, Abner autem et pópulum dormiéntes in circúitu eius. Dixítque Abísai ad David: «Conclúsit Deus hódie inimícum tuum in manus tuas; nunc ergo perfódiam eum láncea in terra semel, et secúndo opus non erit». Et dixit David ad Abísai: «Ne interfícias eum: quis enim exténdit manum suam in christum Dómini et ínnocens erit?». Et dixit David: «Vivit Dóminus quia Dóminus percútiet eum, aut dies eius véniet, ut moriátur, aut in prœ́lium descéndens períbit. Propítius mihi sit Dóminus, ne exténdam manum meam in christum Dómini. Nunc ígitur tolle hastam, quæ est ad caput eius, et scyphum aquæ, et abeámus». Tulit ergo David hastam et scyphum aquæ, qui erat ad caput Saul, et abiérunt; et non erat quisquam, qui vidéret et intellégeret et vigiláret, sed omnes dormiébant, quia sopor Dómini irrúerat super eos.}
\newcommand{\responsoriumi}{\pars{Responsorium 1.} \scriptura{\Rbardot{} 1 Reg. 7, 3 \Vbardot{} ibidem; \textbf{H397}}

\vspace{-5mm}

\responsorium{III}{temporalia/resp-praeparatecordavestra-sinedox.gtex}{}}
\newcommand{\lectioii}{\pars{Lectio II.} \scriptura{1 Sam. 26, 13-25}

\noindent Cumque transísset David ex advérso et stetísset in vértice montis de longe, et esset grande intervállum inter eos, clamávit David ad pópulum et ad Abner fílium Ner dicens: «Nonne respondébis, Abner?». Et respóndens Abner ait: «Quis es tu? Clamásti ad regem!». Et ait David ad Abner: «Numquid non vir tu es? Et quis álius símilis tui in Israel? Quare ergo non custodísti dóminum tuum regem? Ingréssus est enim unus de turba, ut interfíceret regem dóminum tuum. Non est bonum hoc, quod fecísti. Vivit Dóminus quóniam fílii mortis estis vos, qui non custodístis dóminum vestrum, christum Dómini. Nunc ergo vide, ubi sit hasta regis et ubi scyphus aquæ, qui erat ad caput eius».

\noindent Cognóvit autem Saul vocem David et dixit: «Num vox tua hæc est, fili mi David?». Et ait David: «Vox mea, dómine mi rex». Et ait: «Quam ob causam dóminus meus perséquitur servum suum? Quid feci? Aut quod est in manu mea malum? Nunc ergo áudiat, oro, dóminus meus rex verba servi sui: Si Dóminus íncitat te advérsum me, odorétur sacrifícium; si autem fílii hóminum, maledícti sint in conspéctu Dómini, quia eiecérunt me hódie, ut non hábitem in hereditáte Dómini dicéntes: “Vade, servi diis aliénis”. Et nunc non effundátur sanguis meus in terra longe a fácie Dómini; quia egréssus est rex Israel, ut quærat púlicem unum, sicut perséquitur quis perdícem in móntibus».

\noindent Et ait Saul: «Peccávi. Revértere, fili mi David; nequáquam enim ultra malefáciam tibi, eo quod pretiósa fúerit ánima mea in óculis tuis hódie; appáret quod stulte égerim et erráverim multum nimis». Et respóndens David ait: «Ecce hasta regis; tránseat unus de púeris et tollat eam. Dóminus autem retríbuet unicuíque secúndum iustítiam suam et fidem; trádidit enim te Dóminus hódie in manu mea, et nólui exténdere manum meam in christum Dómini. Et sicut magnificáta est ánima tua hódie in óculis meis, sic magnificétur ánima mea in óculis Dómini, et líberet me de omni angústia».

\noindent Ait ergo Saul ad David: «Benedíctus tu, fili mi David; et quidem fáciens fácies et potens póteris». Abiit autem David in viam suam, et Saul revérsus est in locum suum.}
\newcommand{\responsoriumii}{\pars{Responsorium 2.} \scriptura{\Rbardot{} 2 Reg. 7, 8.9.12.13; \Vbardot{} ibid., 7, 9.1; \textbf{H395}}

\vspace{-5mm}

\responsorium{I}{temporalia/resp-egotetuli-CROCHU.gtex}{}}
\newcommand{\lectioiii}{\pars{Lectio III.} \scriptura{Lib. 4, 20, 5-7: SCh 100, 640-642. 644-648}

\noindent Ex Tractátu sancti Irenǽi epíscopi Advérsus hǽreses.

\noindent Vivíficat autem Dei cláritas: percípiunt ergo vitam qui vident Deum. Et propter hoc incapábilis et incomprehensíbilis et invisíbilis visíbilem se et comprehensíbilem et capácem homínibus præstat, ut vivíficet percipiéntes et vidéntes se. Quóniam vívere sine vita impossíbile est, subsisténtia autem vitæ de Dei participatióne évenit, participátio autem Dei est vidére Deum et frui benignitáte eius.

\noindent Hómines ígitur vidébunt Deum ut vivant, per visiónem immortáles facti et pertingéntes usque in Deum. Quod, sicut prædíxi, per prophétas figuráliter manifestabátur quóniam vidébitur Deus ab homínibus qui portant Spíritum eius et semper advéntum eius sústinent. Quemádmodum et in Deuteronómio Móyses ait: \emph{In die ista vidébimus, quóniam loquétur Deus ad hóminem, et vivet.}

\noindent Qui ómnia in ómnibus operátur qualis et quantus est, invisíbilis et inenarrábilis est ómnibus quæ ab eo facta sunt, incógnitus autem nequáquam: ómnia enim per Verbum eius discunt quia est unus Deus Pater, qui cóntinet ómnia et ómnibus esse præstat, quemádmodum in Evangélio scriptum est: \emph{Deum nemo vidit umquam, nisi unigénitus Fílius, qui est in sinu Patris, ipse enarrávit.}

\noindent Enarrátor ergo ab inítio Fílius Patris, quippe qui ab inítio est cum Patre, qui et visiónes prophéticas et divisiónes charísmatum et ministéria sua et Patris glorificatiónem consequénter et compósite osténderit humáno géneri apto témpore ad utilitátem: ubi est enim consequéntia, illic et consonántia; et ubi consonántia, illic et pro témpore; et ubi pro témpore, illic et utílitas.

\noindent Et proptérea Verbum dispensátor patérnæ grátiæ factus est ad utilitátem hóminum, propter quos fecit tantas dispositiónes, homínibus quidem osténdens Deum, Deo autem éxhibens hóminem; et invisibilitátem quidem Patris custódiens, ne quando homo contémptor fíeret Dei et ut semper habéret ad quod profíceret, visíbilem autem rursus homínibus per multas dispositiónes osténdens Deum, ne in totum defíciens a Deo homo cessáret esse: glória enim Dei vivens homo, vita autem hóminis vísio Dei. Si enim quæ est per condiciónem osténsio Dei vitam præstat ómnibus in terra vivéntibus, multo magis ea quæ est per Verbum manifestátio Patris vitam præstat his qui vident Deum.}
\newcommand{\responsoriumiii}{\pars{Responsorium 3.} \scriptura{\textbf{H373}}

\vspace{-5mm}

\responsorium{III}{temporalia/resp-istecognovit-CROCHU-cumdox.gtex}{}

\vfill

\rubrica{vel ad libitum:}

\vspace{3mm}

\pars{Responsorium 3.} \scriptura{\Rbardot{} Ps. 8, 6-8 \Vbardot{} Ps. 20, 4; \textbf{H374}}

\vspace{-5mm}

\responsorium{VII}{temporalia/resp-gloriaethonore-CROCHU-cumdox.gtex}{}}
\newcommand{\hymnuslaudes}{\pars{Hymnus}

\cuminitiali{VI}{temporalia/hym-MartyrDei.gtex}}
\newcommand{\lectiobrevis}{\pars{Lectio Brevis.} \scriptura{2 Cor. 1, 3-5}

\noindent Benedíctus Deus et Pater Dómini nostri Iesu Christi, Pater misericordiárum et Deus totíus consolatiónis, qui consolátur nos in omni tribulatióne nostra, ut possímus et ipsi consolári eos, qui in omni pressúra sunt, per exhortatiónem, qua exhortámur et ipsi a Deo; quóniam, sicut abúndant passiónes Christi in nobis, ita per Christum abúndat et consolátio nostra.}
\newcommand{\responsoriumbreve}{\pars{Responsorium breve.} \scriptura{Ex. 15, 2}

\antiphona{VI}{temporalia/resp-fortitudomea.gtex}}
\newcommand{\preces}{\noindent Fratres, Salvatórem nostrum, testem fidélem,~\gredagger{} per mártyres interféctos propter verbum Dei,~\grestar{} celebrémus, clamántes:

\Rbardot{} Redemísti nos Deo in sánguine tuo.

\noindent Per mártyres tuos, qui líbere mortem in testimónium fídei sunt ampléxi,~\grestar{} da nobis, Dómine, veram spíritus libertátem.

\Rbardot{} Redemísti nos Deo in sánguine tuo.

\noindent Per mártyres tuos, qui fidem usque ad sánguinem sunt conféssi,~\grestar{} da nobis, Dómine, puritátem fideíque constántiam.

\Rbardot{} Redemísti nos Deo in sánguine tuo.

\noindent Per mártyres tuos, qui, sustinéntes crucem,~\gredagger{} tua vestígia sunt secúti,~\grestar{} da nobis, Dómine, ærúmnas vitæ fórtiter sustinére.

\Rbardot{} Redemísti nos Deo in sánguine tuo.

\noindent Per mártyres tuos, qui stolas suas lavérunt in sánguine Agni,~\grestar{} da nobis, Dómine, omnes insídias carnis mundíque devíncere.

\Rbardot{} Redemísti nos Deo in sánguine tuo.}
\newcommand{\benedictus}{\pars{Canticum Zachariæ.} \scriptura{Lc. 12, 8; \textbf{H374}}

\vspace{-4mm}

\antiphona{I f}{temporalia/ant-quimeconfessusfuerit.gtex}

\vspace{-2mm}

\scriptura{Lc. 1, 68-79}

\vspace{-2mm}

\cantusSineNeumas
\initiumpsalmi{temporalia/benedictus-initium-i-f-auto.gtex}

%\vspace{-1.5mm}

\input{temporalia/benedictus-i-f.tex} \Abardot{}}
\newcommand{\benedicamuslaudes}{\cuminitiali{}{temporalia/benedicamus-memoria-laudes.gtex}}
\include{hebdomadaxii}
\include{sabbato}
